\documentclass[UTF8,a4paper,11pt]{ctexart}
\usepackage{geometry}
\geometry{margin=2.2cm}
\usepackage{booktabs}
\usepackage{graphicx}
\usepackage{float}
\usepackage{amsmath,amssymb}
\usepackage{xcolor}
\usepackage{hyperref}
\hypersetup{colorlinks=true,linkcolor=blue,urlcolor=blue}
\title{epsilon-active：精确非光滑 HPWL 与矩形重叠密度全局布局框架}
\author{实验报告}
\date{2026年8月20日}

\begin{document}
\maketitle

\begin{abstract}
本文记录一个独立实现的 C++17/OpenMP 全局布局框架。目标函数严格使用带 pin offset 的
精确 max--min HPWL，密度项严格使用矩形与整数网格 bin 的精确重叠面积；固定宏器件直接
占用 bin 容量。epsilon-active 只改变活动次梯度方向，不改变任何目标值或审计指标。框架
不包含静电模型、目标平滑、合法化、行/站点吸附或合法性检查，并支持多种优化器、最多 40
线程、收敛曲线和布局动画。实验同时报告从 Bookshelf 初始位置的 H24 递归二分负结果，以及
使用外部初始坐标进行精确 recovery 的探索性结果，明确区分两者的证据强度。
\end{abstract}

\section{问题与约束}
挑战 5 的 overflow 用非光滑矩形重叠建模。对 bin $b$，固定宏占用为 $F_b$，可移动器件
的精确占用为 $A_b(x)$，容量为 $C_b=\rho A_b^{\rm bin}-F_b$。报告中的 overflow 是
\[
  \frac{\sum_b \max(0,A_b(x)-C_b)}{\sum_b A_b^{\rm bin}}.
\]
所有矩形--bin 相交面积均由几何裁剪直接计算。固定宏不移动，\texttt{terminal\_NI} 不占用
物理容量。HPWL 对每条 net 直接计算 pin offset 加入后的最大值减最小值；活动集合只用于
选择次梯度方向。没有 WA/LSE、Moreau、静电势、Poisson/DCT 或其它平滑化。

\section{框架实现}
代码位于 \texttt{epsilon-active}。Bookshelf I/O、精确 HPWL oracle、精确 overlap oracle、
epsilon-active、DreamPlace 风格自适应 lambda、Adam/AMSGrad/AdaGrad/Heavy-ball/SGD、精确
回溯、等形状置换和递归容量二分均为独立模块。所有接受的 recovery 移动都重新审计精确
HPWL 与精确 overflow；固定宏永远不进入 movable 集合。OpenMP 线程请求被限制在 40，本文
实验使用本机 16 个逻辑线程。

H24 在固定宏感知容量上递归沿长边切分矩形区域，使用确定性 FM 风格 cut-net refinement，
叶区内按连通性顺序映射到 Hilbert/容量加权锚点。该构造不改变目标定义，但用于检验“先保持
超图连通性再满足容量”的假设。

\section{实验设置}
论文基线 HPWL 取 DREAMPlace 参考值：a1/a2/a3/a4 分别为 73.22M、82.22M、193.72M、
174.08M；允许劣化 1\%，overflow 要求小于 7\%，迭代不超过 5000，运行时间不超过对应
40-thread GP 时间的两倍（67/98/133/187 秒的两倍）。所有本地运行使用 16 线程，因此时间
数字不宣称等价于 40-thread 硬件。

\subsection{H24：clean-start 诊断}
从 Bookshelf 的全零 movable 坐标开始，512x512 bin 的锁定 a1 诊断在 3.65 秒达到
0.5946\% exact overflow，但 HPWL 为 817.219M。叶区 4--64、FM passes 4--20、平衡容差
0.10--0.49 及纳入高阶网络的诊断仍为 747--836M HPWL。H24 的 recovery/swap 延续在
61.91 秒达到 601.391M/6.9998\%，因此没有通过 400M 中间门槛。该负结果说明递归 cut-net
目标和二维锚点映射仍存在结构性失配。

H26 将每个递归 cut 的初始 side 改为当前节点中心所在的物理侧，再执行同一个容量约束 FM
refinement。对非合法化 ePlace GP seed，a1 构造降至 219.7--220.6M HPWL 和 2.1--3.2\%
overflow；20 次 recovery 与 10 次 swap 后为 164.553M/6.8998\%。然而同样设置在 a2/a3/a4
分别为 182.513M、686.471M、627.534M，说明坐标保持只能部分缓解 seed 敏感性。H27 的独立
x/y 容量 rank map 单独留下 62.2553\% overflow，和 H26 组合后为 220.269M/3.2757\%，
因此不具备足够的二维 net 几何保持能力。

H28 将源节点中心和固定宏感知的目标 bin 容量都按同一个整数 Hilbert 次序排序。锁定的 a1
运行在 0.415 秒内达到 0.5409\% exact overflow，但 HPWL 从 41.510M 增至 497.956M，
没有通过 150M 中间门槛。该结果排除了源/目标全局曲线实现不一致，却表明把拥挤 seed 所占的
短曲线区间按全局累计容量拉伸到整张芯片会破坏二维网络几何；单一全局空间填充曲线不足以
同时解决容量和连通性。

H29 进一步把 H26 叶区内的源节点和目标 bin 都改为同一 Hilbert key，但匹配实验得到
220.379M HPWL 和 3.2675\% overflow，略差于 H26 的 219.738M/3.1761\%。因此主要损失
不来自源/目标曲线次序不一致，而来自递归区域分配及一维叶区 quantile 本身。

H30 在每个小叶区内按器件面积降序，并选择离原坐标最近且有固定宏感知剩余容量的二维 bin。
结果的 overflow 改善为 2.0809\%，但 HPWL 为 220.094M，仍未改善 H26。叶区内的容量
规则不是主要瓶颈；节点进入哪个递归物理区域时，长程拓扑损失已经发生。

H31 从同一非合法化 ePlace 坐标直接执行 H8 的八轮精确 excess transport，得到
742.367M HPWL、19.4293\% overflow 和 9.42 秒。相对 H8 全零源坐标的 1,416.36M，
HPWL 损失下降 47.6\%，说明源几何很重要；但逐器件撤离仍会拆散相连节点，未通过 400M
中间门槛。

H32 进一步组合最多 8 个节点的重边组排序，得到 706.182M HPWL、20.1267\% overflow
和 9.323 秒。相对 H31 的 HPWL 仅改善 4.88\%，且 overflow 略微越过 20\% 中间门槛。
原因是当前 group 只改变 contributor 顺序，组员仍独立选择目标；后续需要为连接组联合选择
粗粒度目标容量区域。

H33 让非单节点组的第一个成功目标成为持久二维锚点，后续组员只在锚点邻域内选择严格降低
精确 overflow 的位置。锁定半径 16 得到 635.167M HPWL、18.9867\% overflow 和
10.084 秒，相对 H32 改善 10.06\%，但没有达到预注册的 600.254M 门槛。探索半径
4/8/32 的 HPWL 为 650.910M/645.046M/631.053M；半径 16 到 32 仅改善 4.11M，
时间增加 5.54 秒，因此不继续扩大半径。结果支持共享目标区域，但首个组员仍是较弱的组锚代理。

H34 用全组 incident-net target 的均值替代首个组员 target 进行锚点和局部候选排序，得到
634.210M HPWL、19.0355\% overflow 和 9.919 秒，仅比匹配的 H33 改善 0.15\%，
没有通过 5\% 预注册增益。由此停止 anchor score 调整，后续转向共享网络权重驱动的组成员构造。

H35 按共享网络的 $w/(d-1)$ 权重贪心构造重边组。锁定组大小 8 得到 659.854M HPWL、
19.2340\% overflow 和 9.723 秒，比 H33 反而差 3.89\%；探索组大小 2/4 分别为
657.104M/645.686M，也没有超过 H33。因此 weighted membership 和组大小扫描均被否定。
H29--H35 的外循环复盘发现这些方法共同把器件中心吸附到目标 bin 中心，会清除 ePlace seed
的源 bin 内相对偏移；下一假设将保持源到目标 bin 的共同位移，并继续用精确 overlap oracle
逐候选验收。

H36 对每个候选采用“原中心加目标与源 bin 中心之差”的坐标，从而保持 bin 内偏移。锁定
半径 16 得到 625.958M HPWL、16.8458\% overflow 和 10.140 秒，相对 H33 改善 1.45\%，
但未通过 603.409M 门槛。探索半径 32 为 620.681M/16.6019\%/15.656 秒，相对匹配
H33 也只改善 1.64\%。因此该几何修正确实有效但不是主导损失；下一步需要让连接组共享同一
二维位移，而非仅独立保持各自 bin 内偏移。

H37 在每轮让一个连接组共享首个成功成员的实际二维位移，后续成员越界或不能严格降低精确
overflow 时直接跳过。八轮得到 597.238M HPWL、25.8840\% overflow 和 7.588 秒；
HPWL 相对 H36 改善 4.59\%，但 overflow 未通过 20\% 中间门槛。十二轮为
609.110M/24.4107\%/9.404 秒，说明纯刚性延续的密度收益很弱且损害线长。该结果表明刚性
位移与独立 offset transport 分别偏向拓扑保持和容量疏散，后续应采用预注册的分阶段组合。

H38 固定前八轮刚性、后四轮独立 offset cleanup，结果为 609.939M HPWL、24.2978\%
overflow 和 9.465 秒，未恢复 20\% 门槛。探索 4+8 与 1+11 仍为 24.0942\% 和
23.4760\% overflow，而十二轮全独立对照可达 16.2019\%。因此失败不是 transition 未生效，
而是刚性前缀产生了逐器件下降无法修复的容量碎片。当前组移动仍逐成员要求严格下降；下一结构
方向应把整组位移合并为一次可回滚的精确 overlap 事务，允许单个中间成员为正但组总 delta 为负。

H39 将同一位移的整组矩形变化在临时 occupancy 上顺序求值，并仅按组总 exact overflow
delta 提交或完整回滚。锁定组大小 8 接受 92,305/94,771 次原子尝试，得到 509.126M HPWL、
34.2600\% overflow 和 6.857 秒；HPWL 相对 H37 改善 14.75\%，但密度失败。组大小 4/2
分别为 516.341M/33.3394\% 和 544.112M/30.5289\%，组 2 延长到十六轮仍只有
561.603M/28.2426\%。约 95\% 的高接受率说明 rejected-group splitting 不是主因；后续应
限制整组必须保留首成员一定比例的精确 density gain，否则回滚并仅执行首成员移动。

H40 锁定要求整组保留首成员至少 50\% 的 exact overflow gain，不满足时回滚并进入独立
offset fallback。结果接受 92,170/94,805 组、fallback 2,680 组，得到 522.060M HPWL、
32.6484\% overflow 和 6.943 秒。比例 1.0 端点仍接受 91,460/94,817 组，仅改善到
32.3168\% overflow，因此停止比例调节。根因是原子事务移动全局组的远端非 source 成员并将其
本轮标记完成；下一步应只移动当前 overfull source 中同组的 active contributors。

H41 将原子成员限制为当前 source 的精确矩形 contributors。锁定八轮接受
64,463/65,948 个子组，得到 546.173M HPWL、28.3515\% overflow 和 8.006 秒；相对
H39 改善 5.91 个 overflow 百分点且通过 600M 门槛，但未到 20\%。十六轮仅到
553.817M/27.4672\%，与 ratio-1.0 fallback 组合也只有 556.759M/27.0688\%。因此
source-active atomic 适合作为早期拓扑算子，后续应切换到完全独立的 offset capacity cleanup。

H42 新增可验证的 atomic 轮数前缀：前四轮采用 H41，后八轮进入已有的独立 offset 路径。
锁定运行接受 63,171/63,493 次原子尝试，得到 572.409M HPWL、25.3415\% overflow 和
10.064 秒。HPWL 与时间通过中间门槛，但 overflow 距 20\% 仍差 5.34 个百分点；它也比
十二轮全独立对照高 9.14 个百分点。99.49\% 的接受率说明问题不是原子组经常被拒，而是仅按
总 overflow 下降会允许把较小 excess 搬到新目的 bin，形成后续独立下降无法消除的容量碎片。
因此停止前缀轮数调节，下一步要求原子事务在每个触及 bin 上均不增加精确 overflow。

H43 对每个 atomic trial 进一步要求所有 touched bin 的正 excess 均不增长，否则完整回滚并
进入精确 individual fallback。锁定八轮得到 625.161M HPWL、16.8720\% overflow 和
10.849 秒；密度通过 20\% 门槛，但线长未到 600M。新 guard 拒绝 58,482 次，atomic 仅接受
3,500/62,136 次并 fallback 58,658 次，端点几乎复现 H36 全独立结果。这一结果确认目的 bin
spill 是 atomic 密度 plateau 的原因，也说明只检查首成员选出的单一候选会清除组拓扑收益；
后续应在显式 exact residual capacity 下为 source-active subgroup 搜索并分配多个目的候选。

H44 为每个 source-active subgroup 试算最多八个目的地，并在所有可行项中按组 incident nets
的精确 max-minus-min HPWL delta 选择。锁定运行评估 496,078 个 exact bids，其中 18,308
可行，接受 7,086 组并救回 3,990 个首候选失败组；最终为 622.606M HPWL、16.8107\%
overflow 和 11.486 秒。相对 H43 线长仅改善 2.555M，且未达到 610M 与 10,000 接受组的
预注册门槛。候选可用性并非主导限制；下一步应让同一 source 的全部组在共同 occupancy
snapshot 上生成 bids，再按精确 HPWL 价值竞争并在 live residual capacity 上复验提交。

H45 在每个 source 的共同 occupancy snapshot 上批量生成 bids，按精确 HPWL delta 排序后在
live capacity 上复验。497,424 次试算得到 172,652 个 snapshot-feasible bids，但只提交
3,671/62,411 组；最终 624.363M HPWL、16.8309\% overflow、12.987 秒，比 H44 更差。
这说明大量候选争用同一批稀缺目的容量，单次 greedy priority 不能解决冲突。严格
componentwise-safe group assignment 因此关闭；下一方向应通过同形状 anchor 的精确
density-neutral 身份交换改变 congested source 中需要被疏散的器件，再执行独立 exact transport。

H46 在 transport 前对 source-active connected subgroup 执行同形状身份交换。每个完整交换
均用包含 pin offset 和 net weight 的精确 max-minus-min HPWL 决定接受，并以完整 overlap
oracle 审计密度不变量。锁定运行试算 226,128 个交换，接受 6,683 个、置换 29,560 个器件；
HPWL 从 41.510168M 降至 41.497598M，overflow 在浮点输出上逐位保持
96.8876491423\%。然而该阶段耗时 49.266 秒，随后八轮 H36 最终仅得到
625.990M HPWL、16.7886\% overflow、59.744 秒。密度门槛通过，但 600M 线长与 20 秒
门槛失败。因此关闭局部 connected-group transport/exchange，下一方向改为在 cell-level
cleanup 前显式表示二维全局容量与位移代价。

H47 在 $64\times64$ coarse grid 上建立固定宏容量感知的二维 L1 transport plan，并按完整
incident-net 精确 HPWL delta 实现 offset-preserving 平移。1,871 条 flow edges 和 201,105
次移动在 3.218 秒内把 coarse center-area overflow 从 94.8007\% 降至 0.5257\%。但
$512\times512$ 精确矩形 overlap 仍为 76.2850\%，HPWL 为 530.789M；四轮 H36 后仅到
522.926M/40.8354\%，总时间 8.537 秒。巨大 coarse/exact 差异说明 offset 平移把初始拥塞的
fine-bin phase 复制到了许多 coarse sinks。下一步保留二维全局 flow，但在每个 sink 内用 live
精确矩形容量和精确 HPWL 选择 fine anchors。

H48 对每个 coarse-flow node 枚举其 sink 内全部 64 个 fine-bin centers，并以 live 精确矩形
overflow 严格下降为接受条件。13,433,728 次 anchor 试算提交 200,915 次移动，在 6.804 秒内
把精确 overlap 从 96.8876\% 降至 7.2900\%，验证 H47 的 phase-aliasing 诊断。随后四轮
H36 仅需 5,053 次移动即达到 530.068M HPWL、5.21172\% overflow；总时间 9.595 秒、GP
迭代数 1。H48 的预注册中间门槛全部通过，但 HPWL 仍为 a1 最终上限 73.9522M 的 7.17 倍。
密度表示已不再是主瓶颈；下一步应在每个 source 内联合分配已预留的 coarse sink quotas，按
精确 HPWL opportunity cost 处理匿名 flow 带来的拓扑破坏。

\subsection{H25：外部初始坐标的探索性结果}
H25 使用已有 \texttt{global-placement/ispd2005/*/*.lg.pl} 作为可选初始坐标，然后只运行
本框架的精确 recovery 与等形状置换；程序没有执行合法化或合法性检查。该设置用于隔离
“精确模型是否能保持好种子”这一问题，不能等同于 clean-start 端到端结果。

\begin{table}[H]
\centering
\caption{H25 精确 oracle 结果。时间为 16 线程 wall time，迭代数为 GP 迭代数。}
\begin{tabular}{lrrrrrr}
\toprule
数据集 & HPWL (M) & 基线 (M) & 劣化 & overflow & 时间 (s) & 迭代 \\
\midrule
adaptec1 & 72.7386 & 73.22 & -0.66\% & 6.9000\% & 18.15 & 1 \\
adaptec2 & 81.4895 & 82.22 & -0.89\% & 4.9719\% & 22.07 & 1 \\
adaptec3 & 192.479 & 193.72 & -0.64\% & 4.0422\% & 35.89 & 1 \\
adaptec4 & 174.328 & 174.08 & +0.14\% & 2.3835\% & 34.82 & 1 \\
\bottomrule
\end{tabular}
\end{table}

\section{收敛与可视化}
图中曲线来自 exact recovery/swap CSV，而不是平滑代理。绿色区域表示 overflow 小于 7\%。
每个数据集的布局 GIF 位于 \texttt{output/animations/}，收敛图的 PDF/PNG 位于
\texttt{output/plots/}。

\begin{figure}[H]
\centering
\includegraphics[width=0.92\linewidth]{../output/plots/adaptec1_convergence.pdf}
\caption{adaptec1 的精确 HPWL 与精确 overlap recovery 轨迹。}
\end{figure}
\begin{figure}[H]
\centering
\includegraphics[width=0.92\linewidth]{../output/plots/adaptec2_convergence.pdf}
\caption{adaptec2 的精确轨迹。}
\end{figure}
\begin{figure}[H]
\centering
\includegraphics[width=0.92\linewidth]{../output/plots/adaptec3_convergence.pdf}
\caption{adaptec3 的精确轨迹。}
\end{figure}
\begin{figure}[H]
\centering
\includegraphics[width=0.92\linewidth]{../output/plots/adaptec4_convergence.pdf}
\caption{adaptec4 的精确轨迹。}
\end{figure}

\subsection{H81--H105：严格 \texttt{mu=0} 语义审计}
H81 修正并锁定了同伦语义：overlap 的 epsilon 被强制为零，epsilon-active 只用于
HPWL 的次梯度选择；最后一个 stage 的 \texttt{mu} 被强制写为 $0$，随后仅优化精确
max-minus-min HPWL 与精确矩形 overlap。该严格路径在 H100 首次达到
\(6.24086\%\) exact overflow，但 HPWL 为 316.025M。之后在同一 exact cap 内进行
HPWL 次梯度 recovery、等形状 identity swap 与 incident-net assignment，H105 得到
\(154.050\)M HPWL、\(6.24094\%\) overflow；overflow 在 swap/assignment 中逐位不变。
因此 H105 满足 overflow 门槛，但尚未满足 a1 的 73.9522M HPWL 上限。该差距来自早期
密度构造造成的拓扑损伤，而不是静电项残留或 overlap 平滑。H106 的 position-seeded
recursive bisection 路径得到 170.935M/1.408\%，没有优于 H105。

\begin{table}[H]
\centering
\caption{严格非光滑 clean-start 路径的 a1 审计。}
\begin{tabular}{lrrr}
\toprule
阶段 & HPWL (M) & exact overflow & 说明 \\
\midrule
H81 同伦最终 stage & 62.115 & 82.9201\% & \texttt{mu=0}, overlap epsilon=0 \\
H100 exact breakpoint & 316.025 & 6.24086\% & trust radius 64 \\
H105 feasible recovery & 154.050 & 6.24094\% & exact cap + identity/assignment \\
\bottomrule
\end{tabular}
\end{table}

\section{结论与限制}
精确非光滑模型、固定宏容量和 oracle 保护的 recovery 在高质量初始坐标上满足四个数据集
的数值门槛；这验证了实现的目标/指标一致性和优化器可扩展性。另一方面，H24 clean-start
结果远离论文基线，说明当前容量构造器不能替代一个高质量初始 GP。因而本文不把 H25 的
外部 \texttt{.lg.pl} 锚点结果包装成从零开始的算法胜利。下一步应开发不依赖合法化的非光滑
seed 构造或多尺度超图布局，并在相同的精确 oracle 下重新验证 a1--a4。

\section*{可复现性}
构建：\texttt{mingw32-make -j 16}；测试：\texttt{mingw32-make test}。报告使用指定的
XeLaTeX 2023 编译。核心测试输出为 \texttt{all core tests passed}。

\end{document}
